% ================================================
% 文件: 05_雷达接收机.tex
% 章节: 第10章 雷达接收机
% 所属部分: 接收机技术
% 版本: v1.0
% ================================================
% 本章目录结构（树状）:
% \chapter{雷达接收机}
% ├── \section{雷达接收机概述}
% ├── \section{雷达接收机基本原理}
% ├── \section{雷达接收机主要指标}
% ├── \section{雷达接收机设计要点}
% └── \section{雷达接收机技术发展}
% ================================================

\chapter{雷达接收机}
\label{chap:radar_receiver}

\section{雷达接收机概述}
\label{sec:radar_receiver_intro}

雷达接收机是雷达系统的核心组成部分，
用于接收目标反射的回波信号。
雷达系统通过发射电磁波并接收反射信号来探测目标的距离、
方位和速度。

雷达接收机的主要特点：
\begin{itemize}
    \item \textbf{极高灵敏度}：能够检测微弱的回波信号
    \item \textbf{宽动态范围}：能够处理从微弱回波到强干扰的信号
    \item \textbf{良好的线性度}：确保信号幅度和相位信息准确
    \item \textbf{低噪声}：减少自身噪声对微弱信号的影响
\end{itemize}

雷达接收机与普通通信接收机的最大不同，
在于其输入信号是自身发射脉冲经目标散射后的极微弱回波，
且往往与发射脉冲、
地物杂波、
气象杂波及人为干扰共生。
回波功率随距离的四次方衰减（双程路径损耗），
因此接收机必须具备极低的噪声系数与很高的增益，
同时又要耐得住发射机泄漏与近距离强杂波的冲刷。
现代雷达接收机几乎都采用超外差架构，
并在中频或基带完成脉冲压缩、
杂波抑制与信号检测，
最终把“幅度—距离—速度”信息送给数据处理单元。

从系统角度看，
雷达接收机需与发射机、
天线、
信号处理紧密协同：
它要在发射的极短暂间隙里“开窗”接收，
要承受天线收发开关切换带来的泄漏，
要在强杂波背景中分辨出弱目标。
这些要求决定了雷达接收机在前端保护、
低噪声、
动态范围与线性度上的特殊设计，
后续各节将逐一展开。

雷达按体制可分为脉冲雷达与连续波雷达，
前者以间歇发射的脉冲测距，
后者（如调频连续波 FMCW、
多普勒雷达）靠频率或相位差测距/测速。
按平台与布站又分为单基雷达（收发电离同址）与双基/多基雷达（收发异地）。
工作频段从 VHF、
L、
S、
C、
X 到 Ku、
K、
Ka 乃至毫米波，
频率越高天线越易做小、
分辨越好，
但大气与雨衰减也越强。
接收机须与发射机的脉冲时序、
天线波束严格配合，
构成“发射—空间传播—接收—处理”的完整闭环。

从系统指标看，
雷达接收机追求的是“威力图”上的综合表现：
在给定虚警率下对特定截面积目标达到规定检测概率的最大距离。
这一距离受发射功率、
天线增益、
波长、
目标截面积与接收机最小可检测信号共同决定（雷达方程已给出其关系）。
由于回波功率随距离四次方衰减，
接收机每降低 0.5dB 噪声系数，
往往能换来可观的探测距离增益，
这正是雷达接收机把低噪声置于首位的根本原因。

低截获概率（LPI）是现代雷达的重要设计目标，
而它与接收机灵敏度互为因果：
为降低被敌方侦察到的概率，
雷达压低发射峰值功率、
拉宽 waveform 或采用复杂调制，
这反过来要求自身接收机具备更低的噪声系数与更强的处理增益才能发现目标。
因此 LPI 雷达往往配备数字波束形成、
长时间相参积累与极低相噪本振，
接收机成为达成“看得远又不暴露”的关键瓶颈。

从成本与复杂度看，
接收机（含前端、
A/D 与信号处理）在先进雷达中占比越来越高，
尤其相控阵与数字阵列把“一个接收通道”复制成成百上千个，
幅相一致性、
校准与功耗成为系统级难题。
接收机也不再仅是“收信号”，
而深度参与抗干扰、
自适应波束与模式管理：
它根据环境反馈调整增益、
带宽与处理权重，
是雷达智能化的执行前端。

\section{雷达接收机基本原理}
\label{sec:radar_principle}

雷达接收机的基本工作流程：

1. 
\textbf{接收天线}：
捕获目标反射的回波信号2. 
\textbf{低噪声放大器（LNA）}：
放大微弱信号，
同时保持低噪声3. 
\textbf{混频器}：
将射频信号转换为中频信号4. 
\textbf{中频放大器}：
对中频信号进行高增益放大5. 
\textbf{滤波器}：
滤除干扰和噪声6. 
\textbf{检波器/解调器}：
提取目标信息（距离、
速度、
方位）

雷达方程描述了雷达的探测能力：

\[ R_{max} = \left( \frac{P_t G^2 \lambda^2 \sigma}{(4\pi)^3 S_{min}} \right)^{1/4} \]

其中，$P_t$为发射功率，$G$为天线增益，$\lambda$为波长，$\sigma$为目标截面积，$S_{min}$为最小可检测信号功率。

除雷达方程所揭示的功率链外，
雷达接收机的频域处理同样关键。
回波经天线与环流器进入前端后，
先经过接收机保护器（见设计要点）再送往低噪声放大器。
LNA 之后由混频器与本振将射频搬移到中频，
其关系仍是

\[
f_{\mathrm{IF}}=\left| f_{\mathrm{LO}}-f_{\mathrm{RF}} 
\right|
\]

中频信号经匹配滤波与脉冲压缩后，
幅度随时间（即距离）的分布被提取；
而目标相对平台（如飞机、
舰艇）的径向速度会在回波上叠加多普勒频移

\[
f_d=\frac{2v_r}{\lambda}
\]

其中 $v_r$ 为径向速度，$\lambda$ 为工作波长。借助对一串脉冲的相参处理，即可由多普勒频率分辨出运动目标与静止杂波。

距离分辨由信号带宽决定。对于线性调频（chirp）等脉冲压缩波形，发射的是宽脉冲（以提高平均功率），接收端用匹配滤波器压缩为窄脉冲，压缩后的脉冲宽度约为 $\tau_c\approx 1/B$，对应距离分辨力

\[
\Delta R=\frac{c}{2B}
\]

其中 $c$ 为光速，$B$ 为信号带宽。脉冲压缩使雷达在“大功率远探测”与“高距离分辨”之间得以兼顾，是现代雷达接收机中频/数字处理的核心功能之一。

接收机对回波的处理本质上是一个“匹配滤波 + 二维分辨”问题。匹配滤波器在白噪声下使输出信噪比最大，其对 chirp 信号的响应即为脉冲压缩，输出主瓣宽度约 $1/B$、旁瓣由加权函数决定（如 Hamming、Taylor 加权以抑制距离旁瓣）。距离由回波延迟给出，径向速度由多普勒频率给出，二者构成距离—多普勒二维平面，目标即其中的“亮点”。

脉冲重复频率（PRF）带来一对固有矛盾：PRF 高则无距离模糊、测速好但有距离模糊（远距回波落到近距距离门）；PRF 低则无距离模糊但出现速度模糊（多普勒超出 $\pm f_r/2$ 产生折叠）。工程上常以“高/中/低 PRF”或“多种 PRF 参差”来解距离—速度模糊。此外，chirp 等宽带波形存在距离—多普勒耦合：同一目标因速度引起的多普勒会被误认为距离偏移，需在校正算法中解耦。脉冲压缩、PRF 设计与解模糊共同决定了雷达在密集目标环境下的分辨与识别能力。

距离维的离散化处理依赖距离门或 A/D 采样：模拟接收机按时序选通距离门，数字接收机则对回波采样后按采样点（对应距离量化单元）处理。距离分辨单元 $\Delta R=c/(2f_s)$ 与采样率相关，过粗会丢失近距目标细节。相控阵雷达还引入数字波束形成（DBF）：多路接收通道的数字信号加权求和形成接收波束，可在同一脉冲内同时形成多个波束或自适应置零，接收机因此从“单通道”升级为“多通道阵列处理”。

极化也是接收机可利用的维度：
目标对不同极化（水平/垂直/圆极化）的散射不同，
双极化甚至全极化接收可改善目标识别与杂波抑制。
不过全极化要求通道间极化隔离高、
幅相一致，
进一步抬高前端与校准的复杂度。
这些“多维处理”能力是雷达接收机从传统超外差走向数字阵列后的新特征。

\section{雷达接收机主要指标}
\label{sec:radar_specifications}

雷达接收机的主要技术指标包括：

\begin{itemize}
    \item \textbf{噪声系数（NF）}：衡量接收机产生的噪声大小，通常小于1dB
    \item \textbf{灵敏度}：最小可检测信号功率，通常在-100dBm以下
    \item \textbf{动态范围}：最大与最小可处理信号的比值，通常大于80dB
    \item \textbf{带宽}：接收机能够处理的信号频率范围
    \item \textbf{选择性}：区分目标信号和干扰信号的能力
    \item \textbf{相位噪声}：本地振荡器的相位稳定性
\end{itemize}

噪声系数是雷达接收机最关键的指标之一，
因为它直接决定灵敏度。
整机噪声系数可由各级级联（Friis 公式）求得，
以两级为例：

\[
F=F_1+\frac{F_2-1}{G_1}
\]

其中 $F_1$ 为第一级（LNA）噪声系数，$G_1$ 为其实功率增益，$F_2$ 为第二级噪声系数。式中可见，只要第一级增益足够高、噪声系数足够低，后续混频器、中放的噪声贡献就被显著压制，这正是 LNA 必须置于最前端且追求极低噪声系数的物理原因。优质雷达接收机的整机噪声系数常做到 1\sim3dB，对应灵敏度可达 $-110$ dBm 量级甚至更低。

动态范围描述接收机从最小可检测信号到产生 1dB 增益压缩（或指定杂散）的强信号之间能正常工作的范围，
通常要求大于 80dB。
带宽指标则与波形相关：
中频带宽需与信号（含脉冲压缩后带宽）匹配，
过宽引入多余噪声、
过窄则削掉信号能量。
选择性用于抑制带外杂波与干扰；
相位噪声则关系到本振的短期频率稳定度——本振相位噪声过高会在多普勒处理中抬高杂波基底、
掩盖弱目标，
对脉冲多普勒雷达尤为致命。

检测性能常用两个概率刻画：检测概率 $P_d$ 与虚警概率 $P_{fa}$。虚警指在无目标时误判为有目标，通常要求 $P_{fa}$ 极低（如 $10^{-6}$ 量级）；在给定的信噪比与 $P_{fa}$ 下，接收机可达到相应的 $P_d$。二者共同决定了雷达在杂波与噪声背景中的实际发现能力。

除噪声系数与灵敏度外，
雷达接收机的若干指标相互牵制。
相位噪声决定多普勒滤波的“本底”：
本振近端相位噪声过高会在零多普勒附近抬升杂波基底，
掩埋慢速目标，
对脉冲多普勒雷达尤为关键。
增益压缩点（如 1dB 压缩点）与三阶截获点（IP3）刻画大信号下的线性，
压缩点过低会使强地杂波或副瓣进入时整机增益下降、
弱目标被“淹没”。
中频抑制比与镜频抑制比反映前端滤波对带外与镜频干扰的衰减，
直接影响杂波与干扰下的可用动态。

频率响应的幅度平坦度与群时延一致性也不可忽视：
脉冲压缩对幅相失真敏感，
中频幅频不平或群时延波动会抬高距离旁瓣、
降低分辨。
滤波器矩形度（形状因子）则关系到邻带杂波抑制。
综上，
雷达接收机的指标不是单点最优，
而是在噪声、
线性、
相位噪声、
带宽与选择性之间按任务（搜索/跟踪/成像）统筹的平衡设计；
例如机载下视雷达重杂波环境对相位噪声与动态范围要求极高，
而精密测距雷达更看重带宽与线性。

\section{雷达接收机设计要点}
\label{sec:radar_design}

设计雷达接收机时需考虑：

\begin{itemize}
    \item \textbf{前端设计}：低噪声放大器和混频器的选择
    \item \textbf{频率合成}：稳定的本地振荡器
    \item \textbf{中频处理}：高Q值滤波器和自动增益控制
    \item \textbf{数字信号处理}：AD采样、数字滤波、脉冲压缩
    \item \textbf{抗干扰}：自适应滤波、旁瓣对消
\end{itemize}

\medskip\noindent\textbf{接收机保护器（TR 与限幅）}

雷达发射时，
大功率脉冲经天线泄漏与近距离强反射会进入接收通道，
足以烧毁灵敏的 LNA。
因此前端必须设置接收机保护器：
传统采用放电管（TR 管，
 
transmit-receive tube），
在发射瞬间气体电离形成短路将能量旁路；
现代固态雷达多采用半导体限幅器（PIN 二极管限幅器），
利用大信号下阻抗变化的非线性把尖峰与泄漏功率限制在安全电平。
保护器应有足够低的插入损耗（接收时）与足够高的隔离度（发射时），
并尽量减小对接收机噪声系数的拖累。

\medskip\noindent\textbf{低噪声放大与线性}

LNA 是决定整机噪声系数的瓶颈级，
需在增益、
噪声系数、
输入匹配与线性度之间折中。
混频器则承担变频，
其变频损耗与线性（三阶截获点 IP3）影响动态范围；
为兼顾镜频抑制，
常配合镜像抑制混频或高第一中频方案。
频率合成器须提供低相位噪声、
快速切换的本振，
以保证相参处理与多普勒测量的精度。

\medskip\noindent\textbf{中频/数字处理与脉冲压缩}

中频通道完成放大、滤波后，现代雷达普遍在 A/D 采样后进入数字信号处理：数字下变频、数字滤波、匹配滤波（实现脉冲压缩）、动目标显示（MTI）与脉冲多普勒（PD）处理。MTI 通过对相邻脉冲对消抑制静止杂波，其盲速满足 $v_{\mathrm{blind}}=n\lambda f_r/2$（$f_r$ 为脉冲重复频率）；脉冲多普勒则在多普勒维做滤波，可在强杂波中检测运动目标并测速。抗干扰方面，采用自适应旁瓣对消（ASLC）、旁瓣消隐（SLB）与恒虚警率（CFAR）处理，使检测在起伏杂波中保持稳定的虚警率。

下表归纳雷达接收机设计中的几项核心权衡。

\begin{table}[htbp]
\centering
\caption{雷达接收机设计关键权衡}
\begin{tabular}{|p{4.4cm}|p{5.6cm}|p{5.0cm}|}
\hline
\textbf{设计对象} & \textbf{追求的目标} & \textbf{主要代价/约束} \\
\hline
LNA 噪声系数 & 越低越好以提升灵敏度 & 增益过高压缩动态范围 \\
\hline
接收机保护器 & 高隔离保护前端 & 插入损耗抬升噪声系数 \\
\hline
中频带宽 & 匹配信号保距离分辨 & 过宽引入额外噪声 \\
\hline
本振相位噪声 & 低噪声利多普勒检测 & 频率合成复杂度上升 \\
\hline
脉冲压缩 & 兼顾功率与分辨力 & 需精确匹配滤波与时延 \\
\hline
\end{tabular}
\end{table}

现代雷达接收机越来越多地走向“数字中频/直接采样”。
把 A/D 尽量前移（甚至到射频直接采样），
中频放大、
正交解调、
脉冲压缩、
MTI/PD 与 CFAR 几乎全在数字域完成：
数字下变频得到 I/Q，
匹配滤波用 FFT 或 FIR 实现，
多普勒滤波构成一组窄带滤波器（多普勒滤波器组）。
这带来可重配置、
通道一致好、
标定容易的优点，
代价是对 A/D 采样率、
有效位数（ENOB）与处理算力要求高。

多通道一致性是阵列雷达的生命线：
相控阵每个阵元/子阵一路接收，
通道间增益与相位失配直接决定波束指向与自适应零陷精度，
故需在线或离线通道校准。
抗干扰方面，
自适应旁瓣对消（ASLC）用辅助天线估计干扰、
在主通道相减；
旁瓣消隐（SLB）判别是否来自主瓣；
恒虚警率（CFAR）处理（如单元平均 CFAR）在起伏杂波中自适应设置检测门限，
维持稳定虚警。
上述处理与接收机前端线性、
相位噪声共同决定雷达在电子对抗环境中的生存与探测能力。

收发隔离与保护是接收机可靠工作的前提。
天线收发开关（T/R 开关或环流器）在发射时把大功率导向天线、
阻断接收通道，
接收时反向导通；
固态雷达以 PIN 限幅器作二级保护。
为防止发射机泄漏或近距离强反射烧毁 LNA，
保护链路的隔离度、
响应时间与插入损耗需统筹设计，
任何薄弱环节都会在强发射下暴露。
关键平台常设双接收通道冗余与故障检测，
单一通道失效时无缝切换。

工程实现还强调可测试性与健康管理：
接收机内置测试（BIT）注入已知测试信号，
定期自检噪声系数、
增益与通道一致性，
预报性能退化；
外场则用噪声系数分析仪（Y 因子法）、
矢量信号源与频谱仪实测灵敏度、
动态范围与相位噪声，
与出厂指标比对以判定是否需要返修。
良好的测试与校准体系，
是雷达接收机在全寿命周期内维持探测性能的保障。

\section{雷达接收机技术发展}
\label{sec:radar_development}

雷达接收机技术的发展趋势：

\begin{itemize}
    \item \textbf{数字化}：采用数字接收机技术，提高灵活性和性能
    \item \textbf{软件定义}：软件定义雷达，支持多种工作模式
    \item \textbf{多通道}：多输入多输出（MIMO）技术
    \item \textbf{集成化}：单片微波集成电路（MMIC）和系统级封装（SiP）
    \item \textbf{智能化}：自适应信号处理和机器学习应用
\end{itemize}

数字化是雷达接收机演进的主线。
早期接收机以模拟中频与模拟检波为主，
现代则把 A/D 尽量前移，
形成“数字中频/直接射频采样”接收机，
使脉冲压缩、
MTI/PD、
CFAR 等几乎全在数字域完成，
带来更高的灵活性与可重配置能力。
软件定义雷达（SDR 雷达）在此基础上用同一硬件平台加载不同波形与处理算法，
支持搜索、
跟踪、
成像等多种模式动态切换。

多通道与阵列化是另一重要方向。
相控阵与 MIMO 雷达要求每个阵元或子阵具备独立的接收通道，
通道间幅相一致性直接决定测角与自适应波束形成的性能；
这推动了多通道集成接收机与校准技术发展。
集成化方面，
MMIC 与 SiP 把 LNA、
混频器、
本振、
中放乃至 A/D 集成于微小封装，
显著减小体积、
功耗与成本，
使相控阵的大规模通道成为可能。

智能化则体现在自适应信号处理与机器学习应用：
接收机可根据环境自动调整增益、
带宽、
波形与检测门限，
在复杂电磁对抗中维持探测性能。
上述趋势共同指向“更宽带宽、
更多通道、
更低噪声、
更高集成度与更强算法”的雷达接收机，
以适应隐身目标、
低截获与强对抗的现代需求。

数字阵列雷达（DAR）把收发组件（TR 组件）与数字化下放到阵面，
每个通道独立数字可控，
波束赋形与自适应处理更为灵活；
光控相控阵则用光纤分配本振与中频，
降低馈线损耗与重量，
适合大型阵面。
认知雷达进一步让接收机根据环境反馈自适应选择波形、
频率与处理策略，
在复杂电磁环境中主动规避干扰与杂波。

器件层面，
GaN（氮化镓）功率与低噪器件、
SiGe 与 CMOS 射频收发芯片、
以及系统级封装（SiP）推动接收机向更高集成、
更低功耗发展；
片上雷达（如汽车毫米波雷达单芯片）更把完整收发前端集成于一颗芯片。
未来雷达接收机将在“宽带数字、
多通道一致、
智能处理、
小型集成”四个方向上持续演进，
支撑从远程预警、
精密跟踪到高分辨成像与穿雾探测的广泛任务。

在宽带数字方向上，
直接射频采样与光子模数转换（光子 A/D）有望突破传统电子 A/D 的带宽与位数瓶颈，
使极宽瞬时带宽下的全数字处理成为现实；
多通道一致方向，
片上校准与自补偿算法将把阵列通道失配压到更低，
支撑超大规模阵列。
智能处理方向，
机器学习不仅用于目标识别，
也将用于波形与接收参数自适应调度，
使接收机“按需感知”复杂环境。

小型集成方向，
GaN/SiGe/CMOS 与 SiP 把高功率、
低噪声与数字处理共封装，
推动相控阵与片上雷达走向低成本、
低功耗与可大规模部署；
穿雾、
穿叶与生物医学等特种雷达则对接收机的低噪与线性提出新挑战。
可以预见，
雷达接收机将不再是孤立的“前端盒子”，
而是与天线、
处理、
算法深度耦合的开放式射频感知节点，
持续拓展雷达的能力边界。
